\chapter{Delta-one et dividendes}\label{ch:s2:delta-one-and-dividends}

Une banque qui vend des autocallables se retrouve longue de dividendes qu'elle ne voulait pas, et elle les revend sous forme de contrats à terme sur dividendes et de swaps. Manley et Mueller-Glissmann
ont indiqué qu'à partir de 2004 les swaps de dividendes ont été une source majeure de profit pour les hedge funds multistratégies et macro, car les dividendes implicites de marché se négociaient à
des décotes élevées, alors que les dividendes progressaient. Sur le marché de ce chapitre, les ventes des émetteurs placent les contrats à terme sur dividendes 4~\% sous les dividendes attendus pour chaque année
d'échéance. L'achat du contrat à terme à un an, conservé jusqu'à son échéance, rapporte en moyenne 4.2~\% par an, mais perd une année sur cinq, et de 24~\% lors
d'une année de récession. Le même desk fait de l'\omterm{def:s2:delta-one-and-dividends:indexarb}{arbitrage d'indice} et des \omterm{def:s2:delta-one-and-dividends:financing}{trades de financement} sur le contrat à terme sur indice. La construction est \texttt{firm.deltaone}.

\section{Arbitrage d'indice dans une banque}

\begin{definition}[Arbitrage d'indice]\label{def:s2:delta-one-and-dividends:indexarb}
L'\emph{arbitrage d'indice}\index{arbitrage d'indice} consiste à négocier un contrat à terme sur indice contre le panier de ses composants lorsque le prix du contrat s'écarte de
sa fair value (livre~1, chapitre~21) d'un montant supérieur au coût de négociation des deux jambes, et le débouclage lorsque l'écart se referme.
\end{definition}

\texttt{firm.deltaone} simule une année d'écarts de prix minute par minute du contrat à terme autour de la fair value, avec un écart-type de 4 points de base
de l'indice et une demi-vie de 10 minutes (\cref{lst:s2:delta-one-and-dividends:arb}). Négocier le panier coûte 2 points de base par côté, pour une banque qui
croise en interne une partie de celui-ci (chapitre~25), et 0.5 point de base pour le contrat à terme, soit 5 points de base pour un aller-retour. Le desk vend le contrat à terme et achète le panier lorsque
le contrat à terme est cher de plus de 5 points de base, l'inverse lorsqu'il est bon marché, et procède au débouclage lorsque l'écart se referme~:

\begin{center}
\small
\begin{tabular}{@{}lcccc@{}}
\toprule
délai d'exécution & transactions par an & pb par transaction & pb par an & transactions gagnantes\\
\midrule
aucun & 2\,230 & 1.64 & 3\,657 & 100~\%\\
1 minute & 2\,230 & 1.14 & 2\,539 & 69.2~\%\\
2 minutes & 2\,230 & 0.77 & 1\,707 & 60.7~\%\\
5 minutes & 2\,230 & $-0.26$ & $-584$ & 46.9~\%\\
\bottomrule
\end{tabular}
\end{center}

Sans délai, la transaction ne peut pas perdre, puisqu'elle entre au-delà du coût et sort à zéro. Chaque minute de retard rend une partie d'un écart qui se referme en
dix minutes, et à cinq minutes la transaction perd de l'argent. L'\omterm{def:s2:delta-one-and-dividends:indexarb}{arbitrage d'indice} est une course dont l'enjeu est un point de base, et il revient à celui qui exécute un panier
le plus vite et au moindre coût.

\section{Trades de financement}

\begin{definition}[Trade de financement]\label{def:s2:delta-one-and-dividends:financing}
Un \emph{trade de financement}\index{trade de financement} consiste à détenir un actif face à un dérivé qui le livre plus tard (un panier contre un contrat à terme vendu, ou un
total return swap dans lequel la banque verse à son client le rendement de l'actif, en encaissant le spread de financement implicite dans le prix du dérivé, moins le
coût de financement propre de la banque et de son bilan.
\end{definition}

Le prix d'un contrat à terme au-dessus du prix au comptant, diminué des dividendes, implique un taux auquel le marché finance le panier. Le taux de financement implicite synthétique est le taux overnight
plus un spread qui vaut en moyenne environ 20 points de base et revient à sa moyenne avec une demi-vie de 40 jours, auquel s'ajoutent 25 points de base pendant les dix derniers jours de chaque trimestre,
lorsque les bilans bancaires sont rares (\cref{fig:s2:delta-one-and-dividends:spread}). La banque se finance à 10 points de base au-dessus du taux overnight
et facture 5 au titre du bilan, soit un seuil de 15. L'achat du panier et la vente d'un contrat à terme à trois mois fixent le spread implicite pour le trimestre~:

\begin{center}
\small
\begin{tabular}{@{}p{4.6cm}>{\centering\arraybackslash}p{3.4cm}>{\centering\arraybackslash}p{3.4cm}>{\centering\arraybackslash}p{3.4cm}@{}}
\toprule
entrées chaque trimestre & trimestres retenus & spread moyen verrouillé (pb) & pb de notionnel par an\\
\midrule
au roll & 10~\% & 15.4 & 0.4\\
dans la fenêtre de fin de trimestre & 67.5~\% & 40.7 & 8.0\\
\bottomrule
\end{tabular}
\end{center}

Au roll, le spread franchit rarement le seuil et le coût de l'aller-retour. Dans la fenêtre de fin de trimestre, il le franchit en général~: la banque disposant d'un bilan excédentaire
est rémunérée par ceux qui n'en ont pas. Un total return swap est le même trade sous une autre forme, la banque détenant le panier et versant au client son
rendement contre un spread au-dessus du coût de financement.

\begin{omfigure}
\begin{tikzpicture}
\begin{axis}[width=11cm,height=5.2cm,xlabel={\small année},ylabel={\small spread sur l'overnight (pb)},tick label style={font=\footnotesize},
  xmin=0,xmax=3,ymin=-25,ymax=65,grid=major,grid style={gray!25},table/col sep=comma,
  legend style={font=\scriptsize,at={(0.5,-0.3)},anchor=north},legend columns=2]
\addplot[omThm,thin] table[x=year,y=spread]{figdata/strategies-2/26-delta-one-and-dividends/spread.csv};
\addplot[omDef,thick,dashed] coordinates {(0,15) (3,15)};
\legend{spread de financement implicite du future,seuil de la banque (financement et capital)}
\end{axis}
\end{tikzpicture}

\omcaption{Trois années du spread de financement implicite synthétique du future au-dessus du taux overnight, avec ses sauts de fin de trimestre, face au
seuil de financement et de capital propre à la banque, fixé à 15 points de base. Données~: \texttt{s2\_deltaone.market}.}\label{fig:s2:delta-one-and-dividends:spread}
\end{omfigure}

\section{Risque de dividende issu des émissions structurées}

L'acheteur d'un autocallable est en pratique vendeur d'un put sur l'indice~; l'émetteur détient l'autre côté et couvre son delta en détenant l'indice ou son
contrat à terme. Des dividendes plus élevés abaissent le forward de l'indice, font monter la valeur du put et baissent celle du titre, un gain pour l'émetteur, qui est donc long
de dividendes. Le pricer du livre~5 (chapitre~18) valorise un autocallable à cinq ans avec des observations annuelles, un coupon Phoenix de 8~\% avec une barrière à 70~\% et 60~\%
de protection, à 101.84 pour 100, avec un rendement de dividende de 3~\%. Une hausse de 0.1 point de pourcentage du rendement de dividende la fait baisser de 0.11
(\cref{lst:s2:delta-one-and-dividends:div})~: sur un milliard de titres, soit environ 1.1 million pour 10 points de base de rendement de dividende, sur une durée de vie attendue de
2.66 ans.

\begin{definition}[Pression de l'offre de dividendes]\label{def:s2:delta-one-and-dividends:supply}
La \emph{pression de l'offre de dividendes}\index{pression de l'offre de dividendes} est la décote des contrats à terme ou swaps sur dividendes par rapport aux dividendes attendus qui apparaît lorsque
des hedgers naturellement longs de dividendes, comme les émetteurs de produits structurés, vendent davantage que ce que les autres investisseurs veulent acheter à fair value.
\end{definition}

\begin{dated}{2026-09}{Contrats à terme sur dividendes}\label{dat:s2:delta-one-and-dividends:fexd}
Les EURO STOXX 50 Index Dividend Futures (FEXD) d'Eurex versent 100~EUR par point de dividende de l'indice sur l'année du contrat, avec des échéances allant jusqu'à dix
ans, réglés en espèces le troisième vendredi du mois d'échéance. Eurex décrit les dérivés sur dividendes comme permettant aux investisseurs de prendre des positions sur, ou de se couvrir contre,
les versements futurs de dividendes, notamment pour les produits structurés et les options sur actions.
\end{dated}

\section{Trades sur dividendes}

Le marché synthétique des dividendes affiche des dividendes attendus de 100.8 points d'indice pour l'année suivante, en hausse de 3~\% par an, avec un bruit annuel de 6~\% et une probabilité de 12~\%
chaque année d'une récession qui réduit de 30~\% les dividendes de l'année et l'indice de 25~\%. Les contrats à terme se négocient à des dividendes attendus diminués de 4~\% pour chaque année
d'échéance, la décote de pression de l'offre introduite à dessein, si bien que leur structure par échéance est décroissante alors que les dividendes attendus montent
(\cref{fig:s2:delta-one-and-dividends:dividends}). Un acheteur qui conserve chaque contrat à terme jusqu'à l'échéance, sur 20\,000 scénarios~:

\begin{center}
\small
\begin{tabular}{@{}lccccc@{}}
\toprule
échéance (années) & 1 & 2 & 3 & 4 & 5\\
\midrule
dividendes attendus & 100.8 & 104.7 & 108.7 & 112.9 & 117.2\\
contrat à terme & 96.7 & 96.3 & 95.7 & 94.9 & 93.7\\
rendement moyen par an (\%) & 4.2 & 3.9 & 3.8 & 3.9 & 3.9\\
écart-type annuel (\%) & 12.3 & 9.1 & 7.5 & 6.6 & 5.9\\
quantile à 5~\% par an (\%) & $-25.5$ & $-13.2$ & $-9.4$ & $-8.6$ & $-7.1$\\
probabilité de perte (\%) & 20.9 & 25.2 & 29.8 & 28.0 & 24.6\\
\bottomrule
\end{tabular}
\end{center}

Le rendement de la transaction à un an présente une corrélation de 0.41 avec le rendement de l'indice la même année et perd 24.4~\% en moyenne lors d'une année de récession.
L'acheteur de dividendes est rémunéré pour fournir l'assurance que l'émetteur du produit structuré ne souhaite pas, et il touche ses versements les mêmes années que les actions.
Van Binsbergen, Brandt et Koijen, en reconstituant les prix des STRIPS de dividendes, ont constaté que les créances de dividendes à court terme affichaient des rendements espérés, des
ratios de Sharpe et des volatilités plus élevés que l'indice, avec des bêtas MEDAF inférieurs à un. Kragt, de Jong et Driessen ont ajusté la courbe des contrats à terme sur dividendes de l'EURO STOXX 50 à l'aide de deux
facteurs, l'un revenant à la moyenne en moins d'un an et l'autre sur le cycle économique, et ils ont montré que les investisseurs ne révisent que
faiblement la valeur des dividendes au-delà du cycle économique.

\begin{omfigure}
\begin{tikzpicture}
\begin{axis}[width=11cm,height=5cm,xlabel={\small échéance (années)},ylabel={\small points de dividende de l'indice},tick label style={font=\footnotesize},
  xmin=0.8,xmax=5.2,ymin=90,ymax=120,xtick={1,2,3,4,5},grid=major,grid style={gray!25},table/col sep=comma,
  legend style={font=\scriptsize,at={(0.03,0.97)},anchor=north west}]
\addplot[black,thick,mark=*,mark size=1.2pt] table[x=year,y=expected]{figdata/strategies-2/26-delta-one-and-dividends/dividends.csv};
\addplot[omThm,thick,mark=square*,mark size=1.2pt] table[x=year,y=futures]{figdata/strategies-2/26-delta-one-and-dividends/dividends.csv};
\legend{dividendes attendus,contrats à terme sur dividendes}
\end{axis}
\end{tikzpicture}

\omcaption{Dividendes attendus par année et courbe synthétique des contrats à terme sur dividendes, inférieure de 4~\% par année d'échéance. Données~:
\texttt{s2\_deltaone.dividends}.}\label{fig:s2:delta-one-and-dividends:dividends}
\end{omfigure}

\section{Fiches de stratégie}

\begin{strategyfile}{Arbitrage d'indice}\label{strat:s2:delta-one-and-dividends:arb}
\sfield{Qui vous rémunère, et pourquoi} Des traders qui anticipent le contrat à terme par rapport au panier, ou le panier par rapport au contrat à terme.
\sfield{Instruments et plateformes} Contrats à terme sur indices~; le panier de composantes~; les ETF.
\sfield{Signal} Le contrat à terme comparé à la juste valeur, compte tenu du financement de marché et des dividendes.
\sfield{Dimensionnement et exécution} Au-delà du coût d'un aller-retour~; le panier exécuté en un seul programme.
\sfield{Coûts} Fourchettes du panier et du contrat à terme~; latence.
\sfield{Comment elle meurt} Des concurrents plus rapides~; une entrée erronée de dividende ou de financement.
\sfield{Horizon, capacité, infrastructure} Quelques minutes~; exécution du panier et faible latence.
\sfield{Backtest honnête} Exécuter après le signal, aux prix disponibles à ce moment-là.
\sfield{Sources} Ce chapitre~: 1.14~points de base par transaction avec un délai d'une minute, une perte à cinq minutes.
\end{strategyfile}

\begin{strategyfile}{Financement par total return swap}\label{strat:s2:delta-one-and-dividends:trs}
\sfield{Qui vous rémunère, et pourquoi} Des investisseurs qui veulent le rendement de l'indice sans le financer~; un bilan rare en fin de trimestre.
\sfield{Instruments et plateformes} Total return swaps~; panier contre contrats à terme.
\sfield{Signal} Le spread de financement implicite comparé au coût de financement de la banque et à sa charge en capital.
\sfield{Dimensionnement et exécution} Dans les limites du bilan~; davantage pendant la fenêtre de fin de trimestre.
\sfield{Coûts} Financement~; capital~; négociation du panier.
\sfield{Comment elle meurt} Une réglementation qui relève la charge en capital~; la concurrence pour le bilan.
\sfield{Horizon, capacité, infrastructure} Quelques mois~; trésorerie et allocation du bilan.
\sfield{Backtest honnête} Les coûts de financement et de capital propres à la banque à l'époque.
\sfield{Sources} Ce chapitre~: 8.0~points de base du notionnel par an, saisis en fin de trimestre.
\end{strategyfile}

\begin{strategyfile}{Acheter des dividendes face à l'offre structurée}\label{strat:s2:delta-one-and-dividends:dividends}
\sfield{Qui vous rémunère, et pourquoi} Des émetteurs de produits structurés qui vendent les dividendes que leurs couvertures leur laissent.
\sfield{Instruments et plateformes} Contrats à terme et swaps sur dividendes d'indices.
\sfield{Signal} La décote des contrats à terme par rapport aux dividendes prévus.
\sfield{Dimensionnement et exécution} Conservée jusqu'à l'échéance~; dimensionnée pour une baisse en récession.
\sfield{Coûts} Fourchettes des contrats à terme~; marge.
\sfield{Comment elle meurt} Des baisses de dividendes en récession, y compris les interdictions réglementaires de dividendes bancaires.
\sfield{Horizon, capacité, infrastructure} Un à cinq ans.
\sfield{Backtest honnête} Inclure les années de récession~; prévoir les dividendes à la date de la transaction.
\sfield{Sources} Manley et Mueller-Glissmann (2008)~; van Binsbergen, Brandt et Koijen (2012)~; ce chapitre~: 4.2~\% par an, $-24.4~\%$ pendant une
année de récession.
\end{strategyfile}

\begin{strategyfile}{Calendar spread sur dividendes}\label{strat:s2:delta-one-and-dividends:calendar}
\sfield{Qui vous rémunère, et pourquoi} L'offre concentrée aux échéances dont les produits structurés ont besoin.
\sfield{Instruments et plateformes} Contrats à terme sur dividendes de deux échéances.
\sfield{Signal} La pente de la courbe comparée à la croissance attendue des dividendes.
\sfield{Dimensionnement et exécution} Pondéré de sorte qu'une variation parallèle des dividendes se compense.
\sfield{Coûts} Deux jambes.
\sfield{Comment elle meurt} Une récession qui réduit davantage les dividendes proches que lointains.
\sfield{Horizon, capacité, infrastructure} Des années.
\sfield{Backtest honnête} Un modèle de la structure par terme ajusté sans le contrat à terme.
\sfield{Sources} Kragt, de Jong et Driessen (2020)~; aucun chiffre de performance vérifié.
\end{strategyfile}

\section{Tutoriel~: dividendes dont personne ne voulait}\label{tut:s2:delta-one-and-dividends:tut}

\begin{tutorial}
\textbf{Objectif.} Exécuter l'\omterm{def:s2:delta-one-and-dividends:indexarb}{arbitrage d'indice} avec des délais d'exécution, des \omterm{def:s2:delta-one-and-dividends:financing}{trades de financement} au roll et en fin de trimestre, l'exposition aux dividendes de l'émetteur d'un
autocallable, et un marché de dividendes avec une \omterm{def:s2:delta-one-and-dividends:supply}{pression de l'offre de dividendes}. \textbf{État final~:} les quatre tableaux et les deux figures.
\begin{tutsteps}
\item \textbf{Arbitrage et financement}.
\omcode{code/firm/deltaone/firm_deltaone.py}{78}{104}{Corriger les écarts de prix après un délai~; verrouiller le spread de financement implicite.}\label{lst:s2:delta-one-and-dividends:arb}
\item \textbf{Dividendes}.
\omcode{code/firm/deltaone/firm_deltaone.py}{107}{136}{L'exposition aux dividendes de l'émetteur~; les contrats à terme sur dividendes avec une décote de l'offre.}\label{lst:s2:delta-one-and-dividends:div}
\item \textbf{Exécuter} \texttt{arbitrage()}, \texttt{financing()}, \texttt{exposure()}, \texttt{dividends()} et \texttt{fig\_deltaone.py}.
\end{tutsteps}
\textbf{Pour aller plus loin.} Faire monter la décote de l'offre avec les émissions et la faire baisser après une récession~; couvrir la transaction sur dividendes avec des puts sur indices~; ajouter un
scénario d'interdiction de dividendes pour les banques.
\end{tutorial}

\section{Construction~: delta-one}\label{bld:s2:delta-one-and-dividends:deltaone}

\begin{build}
\bfield{Rôle} \omterm{def:s2:delta-one-and-dividends:indexarb}{Arbitrage d'indice}, \omterm{def:s2:delta-one-and-dividends:financing}{trades de financement}, exposition aux dividendes de l'émetteur et marché de dividendes avec une \omterm{def:s2:delta-one-and-dividends:supply}{pression de l'offre de dividendes}.
\bfield{Interface} \texttt{DeltaOneConfig(\dots)}, \texttt{simulate\_financing(cfg)}, \texttt{index\_arb(sim, cfg, lag)},
\texttt{financing\_trade(sim, cfg, offset)}, \texttt{issuer\_dividend\_exposure(vol, r, q, bump)}, \texttt{dividend\_market(cfg)}.
\bfield{Règles} L'arbitrage au-delà du coût d'un aller-retour, exécuté après un délai~; les \omterm{def:s2:delta-one-and-dividends:financing}{trades de financement} seulement lorsqu'ils franchissent le seuil~; contrats à terme aux
dividendes attendus moins la décote.
\bfield{Tests d'acceptation} \texttt{code/firm/deltaone/tests/}~: pas de délai, pas de perte~; le \omterm{def:s2:delta-one-and-dividends:financing}{trade de financement} à la main~; l'émetteur est acheteur de dividendes~; sans décote,
pas de rendement espéré.
\bfield{Pour aller plus loin} Swaps de dividendes sur actions individuelles~; une offre qui réagit aux émissions~; taux d'intérêt stochastiques.
\end{build}

\begin{omsources}
\item J.~van Binsbergen, M.~Brandt et R.~Koijen, ``On the timing and pricing of dividends'', \emph{American Economic Review} 102(4), 2012.
\item R.~Manley et C.~Mueller-Glissmann, ``The market for dividends and related investment strategies'', \emph{Financial Analysts Journal} 64(3),
2008.
\item J.~Kragt, F.~de Jong et J.~Driessen, ``The dividend term structure'', \emph{Journal of Financial and Quantitative Analysis} 55(3), 2020.
\item Eurex, ``Dividend derivatives'', fiche d'information, consulté le 25 septembre 2026.
\end{omsources}

\section{Exercices}

\begin{exercise}[$\star$]\label{exo:s2:delta-one-and-dividends:1}
Le contrat à terme est surcoté de 7 points de base et un aller-retour coûte 5. L'écart se referme complètement. Que rapporte l'arbitrage sur 50~millions de dollars~?
\end{exercise}

\begin{exercise}[$\star$]\label{exo:s2:delta-one-and-dividends:2}
Le spread implicite vaut 40 points de base, le seuil de la banque à 15, et l'aller-retour d'un trimestre coûte 5 points de base. Que rapporte un \omterm{def:s2:delta-one-and-dividends:financing}{trade de financement}
trimestriel sur 100~millions de dollars~?
\end{exercise}

\begin{exercise}[$\star$]\label{exo:s2:delta-one-and-dividends:3}
Les dividendes attendus valent 108.7 points d'indice et le contrat à terme se négocie à 95.7. Que rapporte sa détention jusqu'à l'échéance si les dividendes arrivent comme prévu, au total et par
an sur trois ans~?
\end{exercise}

\begin{exercise}[$\star\star$]\label{exo:s2:delta-one-and-dividends:4}
Pourquoi l'émetteur d'un autocallable est-il long de dividendes~?
\end{exercise}

\begin{exercise}[$\star\star$]\label{exo:s2:delta-one-and-dividends:5}
Les créances sur dividendes à court terme ont rapporté davantage que l'indice chez van Binsbergen, Brandt et Koijen. Pourquoi l'offre pourrait-elle expliquer une partie de cet écart~?
\end{exercise}

\begin{exercise}[$\star\star$]\label{exo:s2:delta-one-and-dividends:6}
Pourquoi le spread de financement implicite bondit-il en fin de trimestre~?
\end{exercise}

\begin{exercise}[$\star\star\star$]\label{exo:s2:delta-one-and-dividends:7}
\emph{Programmation.} Exécuter \texttt{index\_arb} avec des délais d'exécution de 1, 2 et 5~minutes. À partir de quel délai l'arbitrage cesse-t-il de rapporter, et pourquoi~?
\end{exercise}

\begin{exercise}[$\star\star\star$]\label{exo:s2:delta-one-and-dividends:8}
\emph{Trouver la faille.} «~Les contrats à terme sur dividendes ont rapporté environ 4~\% par an avec une faible volatilité~; ce sont des substituts aux obligations.~»
\end{exercise}

\section{Problème~: Dividendes dont personne ne voulait}

\begin{problem}[{Problème du week-end --- delta-one et dividendes}]\label{pb:s2:delta-one-and-dividends:1}
Le desk synthétique du chapitre et les sources publiques.

\medskip\noindent\textbf{Partie I --- Arbitrage.}
\begin{enumerate}
\item Définir l'\omterm{def:s2:delta-one-and-dividends:indexarb}{arbitrage d'indice}.
\item Décrire l'écart de prix et les coûts.
\item Donner les résultats selon le délai d'exécution.
\item Pourquoi l'arbitrage ne peut-il pas perdre sans délai~?
\end{enumerate}

\medskip\noindent\textbf{Partie II --- Financement.}
\begin{enumerate}[resume]
\item Définir un \omterm{def:s2:delta-one-and-dividends:financing}{trade de financement}.
\item Qu'implique le prix du contrat à terme sur le financement~?
\item Donner les résultats au roll et en fin de trimestre.
\item En quoi un total return swap est-il le même trade~?
\end{enumerate}

\medskip\noindent\textbf{Partie III --- Dividendes.}
\begin{enumerate}[resume]
\item Pourquoi l'émetteur de l'autocallable est-il long de dividendes, et de combien~?
\item Définir la \omterm{def:s2:delta-one-and-dividends:supply}{pression de l'offre de dividendes}.
\item Que verse le contrat Eurex~?
\item Donner le tableau des dividendes.
\end{enumerate}

\medskip\noindent\textbf{Partie IV --- Le verdict.}
\begin{enumerate}[resume]
\item Énoncer le \emph{résultat nommé}~: le rendement du trade sur dividendes issu de la décote liée à l'offre.
\item Que signalent Manley et Mueller-Glissmann~?
\item Qu'ont établi van Binsbergen, Brandt et Koijen~?
\item Qu'ont établi Kragt, de Jong et Driessen~?
\item Quand le trade sur dividendes perd-il de l'argent, et pourquoi à ce moment-là~?
\item Quelle fiche de stratégie exige un bilan~?
\item En quoi ce chapitre se rattache-t-il aux spreads de swap du chapitre~12~?
\item En une phrase~: que vend un desk delta-one~?
\end{enumerate}
\end{problem}

\section{Questions d'entretien}

\begin{interviewq}[$\star$ \iqroles{trader}]\label{iq:s2:delta-one-and-dividends:1}
Comment calculez-vous la fair value d'un contrat à terme sur indice, et quelles données d'entrée peuvent être fausses~?
\end{interviewq}

\begin{interviewq}[$\star\star$ \iqroles{chercheur}]\label{iq:s2:delta-one-and-dividends:2}
Comment prévoiriez-vous les dividendes de l'indice pour l'année prochaine~?
\end{interviewq}

\begin{interviewq}[$\star\star$ \iqroles{trader}]\label{iq:s2:delta-one-and-dividends:3}
Un client souhaite un total return swap de trois mois sur un indice à cheval sur la fin de trimestre. Comment le valorisez-vous~?
\end{interviewq}

\begin{interviewq}[$\star\star$ \iqroles{risque}]\label{iq:s2:delta-one-and-dividends:4}
Quel test de résistance appliqueriez-vous à un portefeuille long de contrats à terme sur dividendes~?
\end{interviewq}

\begin{interviewq}[$\star\star$ \iqroles{développeur}]\label{iq:s2:delta-one-and-dividends:5}
Concevez le système qui maintient les fair values des indices à jour en cours de séance.
\end{interviewq}

\begin{interviewq}[$\star\star\star$ \iqroles{chercheur}]\label{iq:s2:delta-one-and-dividends:6}
Le prix du contrat à terme est $F = (S - D) e^{rT}$ où $D$ est la valeur actuelle des dividendes jusqu'à l'échéance. Montrez qu'une erreur $\delta$ sur $D$ décale la fair value de
$-\delta e^{rT}$, et expliquer pourquoi une erreur de prévision sur les dividendes ressemble à un arbitrage.
\end{interviewq}
